\let\R\relax
\newcommand{\R}{\mathbb{R}}
\let\C\relax
\newcommand{\C}{\mathbb{C}}
\let\N\relax
\newcommand{\N}{\mathbb{N}}
\let\Z\relax
\newcommand{\Z}{\mathbb{Z}}
\let\Q\relax
\newcommand{\Q}{\mathbb{Q}}
\let\cf\relax
\newcommand{\cf}{\operatorname{cf}}
\let\pp\relax
\newcommand{\pp}{\operatorname{pp}}
\let\tcf\relax
\newcommand{\tcf}{\operatorname{tcf}}
\let\otp\relax
\newcommand{\otp}{\operatorname{otp}}
\let\add\relax
\newcommand{\add}{\operatorname{add}}
\let\dom\relax
\newcommand{\dom}{\operatorname{dom}}
\let\ran\relax
\newcommand{\ran}{\operatorname{ran}}
\let\Hull\relax
\newcommand{\Hull}{\operatorname{Hull}}
\let\Skolem\relax
\newcommand{\Skolem}{\operatorname{Skolem}}
\let\DC\relax
\newcommand{\DC}{\operatorname{DC}}
\let\MA\relax
\newcommand{\MA}{\operatorname{MA}}
\let\CH\relax
\newcommand{\CH}{\operatorname{CH}}
\let\GCH\relax
\newcommand{\GCH}{\operatorname{GCH}}
\let\ZFC\relax
\newcommand{\ZFC}{\operatorname{ZFC}}
\let\ZF\relax
\newcommand{\ZF}{\operatorname{ZF}}
\let\Spec\relax
\newcommand{\Spec}{\operatorname{Spec}}
\let\pcf\relax
\newcommand{\pcf}{\operatorname{pcf}}
\let\Ext\relax
\newcommand{\Ext}{\operatorname{Ext}}

\chapter{Saharon Shelah (2001)}
\index{Shelah, Saharon}

\begin{quote}
\it ``Shelah has reshaped model theory, turning it from a branch of logic into a powerful tool with applications throughout set theory, algebra, and analysis. His classification theory provided the first deep structure theorem for mathematical theories, and his pcf theory revolutionized cardinal arithmetic. He is one of the most prolific mathematicians in history, with over 1200 publications.'' --- Wolf Prize Citation
\end{quote}

Saharon Shelah (born July 3, 1945) is one of the most prolific and influential mathematicians in history. The 2001 Wolf Prize in Mathematics (shared with Vladimir Arnold) recognized his foundational contributions to model theory, set theory, and the foundations of mathematics. With over 1200 published papers and numerous books, Shelah's work has transformed the landscape of mathematical logic and extended its reach into algebra, topology, and analysis.

Shelah's mathematical style is characterized by extraordinary technical power, a relentless drive for generality, and a knack for formulating the right definitions to capture deep structural phenomena. His classification theory brought order to the chaos of first-order theories, revealing that every theory falls into one of a few stable classes with profound structural consequences. His pcf (possible cofinalities) theory shattered the long-standing impasse in cardinal arithmetic, proving deep theorems about the behavior of the continuum function on singular cardinals. His development of proper forcing provided a systematic framework for independence proofs. His solution of the Whitehead problem --- showing it is independent of ZFC --- was a landmark of set-theoretic algebra. And his work on J\'onsson algebras, saturated models, and the Keisler--Shelah theorem on ultraproducts has had lasting impact on model theory.

\section{Early Life and Career}

Saharon Shelah was born on July 3, 1945, in Jerusalem, then in the British Mandate for Palestine (now Israel). His father, Yonatan Shelah, was a poet and translator, and his mother, Batya, was a teacher. Shelah showed extraordinary mathematical talent from an early age, excelling in mathematics and physics during his secondary education.

Shelah entered the Hebrew University of Jerusalem in 1962, where he studied mathematics and physics. He completed his Bachelor of Science in 1964 and his Master of Science in 1967, both at the Hebrew University. His early interests included both model theory and set theory, and he was influenced by the work of Abraham Robinson, Alfred Tarski, and other pioneers of mathematical logic.

Shelah completed his PhD at the Hebrew University of Jerusalem in 1969 under the supervision of Michael Rabin. His doctoral dissertation, ``Stable Theories,'' laid the groundwork for what would become his monumental classification theory. In this work, he introduced and systematically developed the concept of stability for first-order theories, proving the first deep structure theorems about the models of stable theories.

After completing his PhD, Shelah spent a year as a postdoctoral fellow at the University of California, Berkeley (1969--1970), where he worked with Robert Vaught and others. He then joined the faculty of the Hebrew University of Jerusalem, where he has remained for most of his career. He also holds a long-term appointment at Rutgers University, where he spends part of each year.

The 1970s were an extraordinarily productive period for Shelah. During this decade, he:
\begin{itemize}
    \item Developed the foundations of classification theory, including the stability hierarchy and the structure--nonstructure theorem.
    \item Solved the Whitehead problem in abelian group theory, showing it is independent of ZFC.
    \item Began his systematic study of forcing, developing proper forcing and later semiproper forcing.
    \item Made fundamental contributions to the theory of saturated models and ultraproducts.
\end{itemize}

In the 1980s, Shelah turned his attention to cardinal arithmetic, developing pcf (possible cofinalities) theory. This work revolutionized the field, providing new tools for understanding the behavior of the continuum function on singular cardinals and proving results that had been considered out of reach.

The 1990s saw Shelah extend his work on forcing to even more general frameworks, including semiproper forcing, iterated forcing, and the analysis of the consistency strength of various set-theoretic statements. He continued to produce deep results in model theory, including the final classification of the stable theories and the analysis of various non-elementary frameworks.

Shelah was awarded the Wolf Prize in Mathematics in 2001, shared with Vladimir Arnold. The prize citation emphasized his ``fundamental contributions to mathematical logic, model theory, and set theory, and their applications in other branches of mathematics.'' He has also received the Israel Prize (1998), the EMET Prize (2011), and the Rothschild Prize (1982). He is a member of the Israel Academy of Sciences and Humanities, the American Academy of Arts and Sciences, and the National Academy of Sciences of the United States.

\section{Classification Theory}

Classification theory is Shelah's most systematic and far-reaching contribution to mathematics. It addresses a fundamental question: given a first-order theory \(T\), what can we say about the structure of its models? Can we classify the models of \(T\) up to isomorphism by some reasonable set of invariants?

Prior to Shelah's work, the study of complete first-order theories had produced isolated results but no systematic framework. The L\"owenheim--Skolem theorem, the Compactness theorem, and G\"odel's completeness theorem had established the basic properties of first-order logic, and the work of Vaught, Morley, and others had begun the analysis of countable models. But the general question of what makes a theory ``tame'' or ``wild'' remained mysterious.

Shelah's key insight was to introduce the notion of \emph{stability} as a dividing line. A theory \(T\) is called \(\kappa\)-stable if for every model \(M\) of \(T\) and every subset \(A \subseteq M\) of size at most \(\kappa\), the set of complete types over \(A\) has size at most \(\kappa\). If \(T\) is \(\kappa\)-stable for some infinite \(\kappa\), it is said to be \emph{stable}.

\begin{definition}
Let \(T\) be a complete first-order theory in a countable language. For an infinite cardinal \(\kappa\), \(T\) is \textbf{\(\kappa\)-stable} if for every model \(M \models T\) and every \(A \subseteq M\) with \(|A| \leq \kappa\), the set \(S(A)\) of complete types over \(A\) (relative to \(T\)) has cardinality at most \(\kappa\). The theory \(T\) is called \textbf{stable} if it is \(\kappa\)-stable for some infinite cardinal \(\kappa\).
\end{definition}

\begin{theorem}[Shelah's Stability Spectrum Theorem, 1971]
For any complete countable first-order theory \(T\), the set of cardinals \(\kappa\) for which \(T\) is \(\kappa\)-stable is either:
\begin{itemize}
    \item Empty (the theory is \emph{unstable}),
    \item All infinite cardinals (the theory is \emph{superstable} and every cardinal is stable),
    \item All cardinals \(\kappa\) with \(\kappa \geq \beth_{n}(|\cdot|)\) for some finite \(n\) (the theory is \emph{strictly stable}),
    \item All cardinals \(\kappa \geq \beth_{\omega}(|\cdot|)\) (the theory is \emph{strictly superstable}).
\end{itemize}
Moreover, there is a deep connection between the stability spectrum and the existence of certain combinatorial configurations in the models of \(T\).
\end{theorem}

The classification of theories by stability refines the earlier work of Morley, who had shown that a countable theory that is \(\kappa\)-stable for some uncountable \(\kappa\) must be \(\kappa\)-stable for all \(\kappa\). Shelah's spectrum theorem provides a complete description of the possible stability spectra.

The real power of stability theory is captured by the \emph{structure--nonstructure theorem}. Shelah showed that stable theories admit a rich structure theory: their models can be analyzed using notions of independence, forking, and regularity, and the class of models of a stable theory can be classified up to isomorphism by a set of invariants (the \emph{Shelah--Lascar classification}). Unstable theories, by contrast, are ``wild'': they have the maximal number of models in all uncountable cardinalities, and their model theory admits no reasonable classification.

\begin{theorem}[Shelah's Structure--Nonstructure Theorem, 1970s]
Let \(T\) be a complete countable first-order theory. Either:
\begin{enumerate}
    \item \(T\) is stable, in which case the models of \(T\) admit a structure theory based on forking and independence, and the isomorphism types of models of \(T\) in each uncountable cardinal can be described via dimension-like invariants; or
    \item \(T\) is unstable, in which case for every uncountable cardinal \(\kappa\), \(T\) has \(2^{\kappa}\) non-isomorphic models of cardinality \(\kappa\) (the maximal possible number).
\end{enumerate}
\end{theorem}

The structure--nonstructure theorem provides a rigorous mathematical formulation of the idea that stable theories are ``tame'' and unstable theories are ``wild.'' The dividing line is precisely the existence of the order property: a theory is unstable if and only if it has a formula that defines a partial order with an infinite chain (the \emph{order property}).

Shelah's classification theory has had a profound impact beyond model theory. In algebra, it showed that the theory of algebraically closed fields is stable (indeed, uncountably categorical), and the classification of its models by transcendental degree is a special case of the general structure theory. The theory of differentially closed fields, separably closed fields, and other algebraic structures have been analyzed through the lens of stability theory. In analysis, the theory of compact complex manifolds has been shown to be stable in certain cases, leading to new insights into the geometry of complex structures.

\subsection{The Stability Hierarchy}

Shelah refined the classification of stable theories by introducing finer dividing lines. A theory is:
\begin{itemize}
    \item \textbf{Unstable} if it has the order property: there is a formula \(\varphi(x,y)\) such that in some model, \(\varphi\) defines an infinite linear order.
    \item \textbf{Strictly stable} if it is stable but not superstable. These theories have a complicated behavior with respect to the ranks of types.
    \item \textbf{Superstable} if it is stable and every type has ordinal-valued rank. Superstable theories are closer to the categorical case and admit a cleaner structure theory.
    \item \textbf{\(\omega\)-stable} if it is \(\omega\)-stable (also called totally transcendental). These are the tamest theories, generalizing the properties of algebraically closed fields.
\end{itemize}

\begin{theorem}[Shelah, 1970]
A countable complete first-order theory \(T\) is \(\omega\)-stable if and only if it has no formula that defines an infinite set of pairwise-inconsistent formulas (no \emph{strict order property}). Equivalently, \(T\) is \(\omega\)-stable if and only if the Morley rank of every type is defined and ordinal-valued.
\end{theorem}

The stability hierarchy is not merely a classification for its own sake. Each level of the hierarchy corresponds to a specific kind of combinatorial behavior that can be exploited or, in the unstable case, that prevents classification.

\subsection{Forking and Independence}

The central technical concept in stability theory is \emph{forking}, introduced by Shelah as a model-theoretic analogue of algebraic independence. A formula \(\varphi(x,b)\) is said to \emph{fork} over a set \(A\) if it implies a certain kind of dependence on the parameter \(b\) that cannot be eliminated. In a stable theory, forking satisfies the following fundamental properties:

\begin{definition}[Forking]
Let \(T\) be a stable theory. A formula \(\varphi(x,b)\) \textbf{forks} over a set \(A\) if there exist \(b_0,b_1,\dots\) and automorphisms \(\sigma_i\) fixing \(A\) such that \(\{\sigma_i(\varphi(x,b)) : i < \omega\}\) is \(k\)-inconsistent for some finite \(k\). Equivalently, \(\varphi(x,b)\) does not fork over \(A\) if it can be extended to a complete type over any superset of \(A\).
\end{definition}

Shelah proved that in a stable theory, forking satisfies:
\begin{enumerate}
    \item \textbf{Invariance:} Forking is invariant under automorphisms.
    \item \textbf{Monotonicity:} If \(\varphi\) forks over \(A\) and \(A \subseteq B\), then \(\varphi\) forks over \(B\).
    \item \textbf{Transitivity:} Forking over a set and then over a larger set is equivalent to forking over the larger set.
    \item \textbf{Local character:} Every type does not fork over some subset of size at most \(|\!T\!|\).
    \item \textbf{Stationarity:} Over a model, types are stationary: if two types are non-forking extensions of the same type over a model, they are equal.
    \item \textbf{The Independence Theorem:} If \(A\) and \(B\) are independent over \(C\) (in the sense of forking), then any two types over \(A\) and \(B\) that agree on \(C\) can be extended to a complete type over \(A \cup B\).
\end{enumerate}

These properties make forking the correct model-theoretic analogue of linear independence in vector spaces, algebraic independence in fields, and free amalgamation in combinatorial structures. The existence of a notion of independence satisfying these properties is what makes stable theories tractable.

\section{The Whitehead Problem}

One of Shelah's most celebrated results is his solution of the Whitehead problem in abelian group theory. The problem, posed by J. H. C. Whitehead in the 1950s, is deceptively simple:

\begin{condition}[Whitehead's Problem]
Let \(A\) be an abelian group. Is every extension of \(\Z\) by \(A\) necessarily split? Equivalently, does \(\Ext^1(A,\Z) = 0\) imply that \(A\) is free? Such groups are called \emph{Whitehead groups} or \emph{W-groups}.
\end{condition}

More precisely, \(\Ext^1(A,\Z)\) classifies extensions \(0 \to \Z \to G \to A \to 0\) up to equivalence. The group \(A\) is a Whitehead group (or W-group) if \(\Ext^1(A,\Z) = 0\). Every free abelian group is trivially a Whitehead group. The question is whether the converse holds: are all Whitehead groups free?

The problem was known to be difficult. It was studied by Stein, Fuchs, Chase, and others, who proved various partial results. For countable groups, Stein (1951) proved that every countable Whitehead group is free. But the uncountable case remained wide open.

Shelah's resolution of the problem is one of the landmarks of set-theoretic algebra. In 1974, he proved that the Whitehead problem is independent of ZFC: it can neither be proved nor disproved from the standard axioms of set theory.

\begin{theorem}[Shelah's Independence Theorem, 1974]
\begin{enumerate}
    \item If \(V = L\) (G\"odel's axiom of constructibility), then every Whitehead group is free. Equivalently, \(\Ext^1(A,\Z) = 0\) implies \(A\) is free.
    \item If \(\MA + \neg\CH\) (Martin's Axiom plus the negation of the Continuum Hypothesis) holds, then there exists a non-free Whitehead group of cardinality \(\aleph_1\).
\end{enumerate}
Thus the statement ``every Whitehead group is free'' is independent of ZFC.
\end{theorem}

The proof of the first statement uses the combinatorial principles that hold in \(L\), particularly the existence of a \(\diamondsuit\)-sequence, to construct an unsplit extension for any non-free group. The proof of the second statement uses Martin's Axiom to construct a non-free group all of whose countable subgroups are free, and then to show that such a group must be a Whitehead group.

The Whitehead problem was the first purely algebraic problem shown to be independent of ZFC. It demonstrated that even apparently concrete questions about abelian groups can transcend the power of the standard set-theoretic axioms. This result had a profound impact on both algebra and set theory, leading to the systematic study of the set-theoretic content of algebraic statements and the development of new forcing techniques.

\begin{remark}
Shelah's result can be seen as establishing that the Whitehead problem is not about algebra alone: the answer depends on the set-theoretic universe in which the algebra is done. This was a revelation for many mathematicians, showing that independence phenomena are not confined to the rarefied heights of large cardinals and the continuum hypothesis but can manifest in concrete algebraic questions.
\end{remark}

The Whitehead problem led to an extensive literature on set-theoretic algebra. Shelah himself proved many related results: for example, the existence of a non-free Whitehead group of cardinality \(\aleph_n\) for any finite \(n\) is also independent of ZFC, but the question for cardinality \(\aleph_\omega\) is provably false in ZFC (every Whitehead group of cardinality \(\aleph_\omega\) is free). This subtle interplay between set theory and algebra is characteristic of Shelah's work.

\section{PCF Theory}

Perhaps Shelah's most stunning technical achievement is pcf (possible cofinalities) theory, developed in the 1980s and presented in his monumental book \textit{Cardinal Arithmetic} (1994). PCF theory revolutionized the study of cardinal arithmetic, which had been largely stagnant since the work of G\"odel, Cohen, and Easton in the 1960s.

The fundamental problem of cardinal arithmetic is to determine the values of \(\aleph_\alpha^{\aleph_\beta}\) as a function of \(\aleph_\alpha\) and \(\aleph_\beta\). The classical theorems of K\"onig, Hausdorff, and Tarski had established basic inequalities:
\[
\aleph_{\alpha+1} \leq \aleph_\alpha^{\aleph_\alpha} \leq \aleph_{\alpha+1}^{\aleph_\alpha},
\]
and K\"onig's theorem gives:
\[
\cf(\aleph_\alpha^{\aleph_\beta}) > \aleph_\beta.
\]
But beyond these elementary facts, very little was known about the exact values of cardinal exponentials, especially at singular cardinals.

The problem of singular cardinals is particularly acute. Easton's theorem (1970) shows that for regular cardinals, the continuum function \(\kappa \mapsto 2^\kappa\) can be arbitrarily large subject only to the constraints of monotonicity and K\"onig's theorem. But for singular cardinals, the behavior is much more constrained, and the classical results gave very little information.

Shelah's pcf theory provided a powerful new framework for analyzing the arithmetic of singular cardinals. The central concept is the set of \emph{possible cofinalities} of a reduced product of regular cardinals.

\begin{definition}[PCF]
Let \(\mathfrak{a}\) be a set of regular cardinals. For any ultrafilter \(D\) on \(\mathfrak{a}\), the ultrapower \(\prod \mathfrak{a} / D\) is a linearly ordered set; its cofinality is a regular cardinal called a \textbf{possible cofinality} of \(\mathfrak{a}\). The set of all possible cofinalities of \(\mathfrak{a}\) is denoted \(\pcf(\mathfrak{a})\):
\[
\pcf(\mathfrak{a}) = \{\cf(\prod \mathfrak{a} / D) : D \text{ is an ultrafilter on } \mathfrak{a}\}.
\]
\end{definition}

\begin{theorem}[Shelah's PCF Theorem, 1980s]
Let \(\mathfrak{a}\) be a set of regular cardinals with \(|\mathfrak{a}| < \min(\mathfrak{a})\). Then:
\begin{enumerate}
    \item \(\pcf(\mathfrak{a})\) has a maximum element, denoted \(\max \pcf(\mathfrak{a})\).
    \item \(\max \pcf(\mathfrak{a}) \leq |\mathfrak{a}|^{+4}\), and more generally, \(\pcf(\mathfrak{a})\) has cardinality at most \(2^{|\mathfrak{a}|}\).
    \item If \(\lambda = \max \pcf(\mathfrak{a})\), then the cofinality of \(\lambda\) is equal to the cofinality of \(\prod \mathfrak{a} / D\) for some ultrafilter \(D\).
    \item There exists a sequence of \(\lambda\)-complete ultrafilters that generate all elements of \(\pcf(\mathfrak{a})\).
\end{enumerate}
\end{theorem}

The most famous corollary of pcf theory is Shelah's inequality:

\begin{theorem}[Shelah's Inequality, 1980]
For every infinite cardinal \(\aleph_\omega\),
\[
\aleph_\omega^{\aleph_0} < \aleph_{\omega_4}.
\]
More generally, if \(\aleph_\omega\) is a strong limit cardinal, then \(2^{\aleph_\omega} < \aleph_{\omega_4}\).
\end{theorem}

This result was shocking to the mathematical community. Before Shelah's work, it was consistent with the known axioms of set theory that \(2^{\aleph_\omega}\) could be arbitrarily large (subject to the cofinality constraint that \(\cf(2^{\aleph_\omega}) > \aleph_0\)). Shelah's theorem showed that, contrary to expectations, \(2^{\aleph_\omega}\) is bounded by \(\aleph_{\omega_4}\). The bound is not absolute --- it depends on the underlying cardinal arithmetic --- but it provides an explicit and relatively low upper bound that holds in every model of ZFC.

The number \(\omega_4\) in the theorem has a special significance. Shelah's original proof gave a bound of \(\aleph_{(2^{\aleph_0})^{+}}\), which was later improved through refinements of the pcf theory to \(\aleph_{\omega_4}\). It is known that \(\aleph_{\omega_4}\) is the optimal upper bound in the sense that it is consistent that \(\aleph_\omega^{\aleph_0} = \aleph_{\omega_1}\), but it is an open problem whether the bound can be improved to \(\aleph_{\omega_3}\) or \(\aleph_{\omega_2}\).

\subsection{The Continuum Function on Singular Cardinals}

One of the most powerful applications of pcf theory is to the analysis of the continuum function \(\kappa \mapsto 2^\kappa\) on singular cardinals. The classical result of Silver (1974) showed that if the continuum function is constant below a singular cardinal of uncountable cofinality, then it is also constant at that singular cardinal. Shelah extended Silver's theorem to a much broader context.

\begin{theorem}[Shelah's Singular Cardinal Bound, 1980s]
Let \(\kappa\) be a singular cardinal. Suppose that for a set of cardinals below \(\kappa\) that is cofinal in \(\kappa\), the continuum function takes values that are not too large. Then \(2^\kappa\) is bounded above by a function of the cofinality of \(\kappa\) and the values of the continuum function below \(\kappa\).

More specifically, for any singular cardinal \(\kappa\),
\[
2^\kappa < \aleph_{(2^{\cf(\kappa)})^{+}}.
\]
\end{theorem}

This result shows that the continuum function on singular cardinals is severely constrained, in sharp contrast to the situation for regular cardinals where Easton's theorem shows it can be arbitrarily large.

\subsection{The \texorpdfstring{\(\pp(\kappa)\)}{pp(kappa)} Function}

Shelah introduced the function \(\pp(\kappa)\) (``possible powers''), defined as the supremum of all cardinals of the form \(\cf(\prod \mathfrak{a} / D)\) where \(\mathfrak{a}\) is a set of regular cardinals with \(\sup \mathfrak{a} = \kappa\) and \(D\) is an ultrafilter on \(\mathfrak{a}\). The function \(\pp(\kappa)\) captures the essential behavior of cardinal exponentials at \(\kappa\).

\begin{theorem}[Properties of \(\pp(\kappa)\)]
For any cardinal \(\kappa\):
\begin{enumerate}
    \item \(\kappa^{+} \leq \pp(\kappa) \leq 2^{\kappa}\).
    \item If \(\kappa\) is regular, \(\pp(\kappa) = \kappa^{+}\).
    \item If \(\kappa\) is singular, \(\pp(\kappa)\) can be larger than \(\kappa^{+}\), but is bounded by \(2^{\kappa}\).
    \item \(\pp(\kappa)\) is regular.
    \item \(\pp(\kappa)\) depends only on the cofinality of \(\kappa\) and the sequence of cardinals below \(\kappa\) that are relevant for the pcf analysis.
\end{enumerate}
\end{theorem}

The \(\pp\) function has been extensively studied and is now understood to be the fundamental object of cardinal arithmetic. Many classical results, including Silver's theorem, can be elegantly reformulated in terms of \(\pp\).

Shelah's work on cardinal arithmetic, culminating in his book \textit{Cardinal Arithmetic} (1994), is one of the most impressive technical achievements in set theory. The pcf theory has found applications far beyond its original context, including in the study of the singular cardinals hypothesis, the analysis of ultrafilters and ultrapowers, and even in combinatorial number theory.

\section{Forcing and Independence Results}

Shelah has been one of the most important contributors to the theory of forcing, the technique invented by Paul Cohen in 1963 for proving independence results in set theory. Shelah's work on forcing has transformed it from a collection of ad-hoc constructions into a systematic theory with powerful general frameworks.

\subsection{Proper Forcing}

The notion of \emph{proper forcing}\index{proper forcing axiom}, introduced by Shelah in the 1970s, is one of the most important concepts in modern set theory. Proper forcing was motivated by the need to understand which forcing notions preserve certain properties of the ground model, particularly the stationarity of subsets of \(\omega_1\).

\begin{definition}[Proper Forcing]
Let \(\mathbb{P}\) be a forcing notion. \(\mathbb{P}\) is \textbf{proper} if for every countable elementary submodel \(N \prec H(\theta)\) for some sufficiently large \(\theta\) with \(\mathbb{P} \in N\), and for every condition \(p \in \mathbb{P} \cap N\), there exists a condition \(q \leq p\) that is \(N\)-generic: for every dense open set \(D \subseteq \mathbb{P}\) with \(D \in N\), the set \(q \cap D \cap N\) is non-empty.
\end{definition}

The concept of proper forcing captures the idea that the forcing preserves properties of the ground model in a strong sense. Proper forcing notions include many of the most important constructions in set theory: ccc (countable chain condition) forcing, closed forcing, and many others.

\begin{theorem}[Shelah's Proper Forcing Theorem, 1978]
Proper forcing notions have the following properties:
\begin{enumerate}
    \item Proper forcing preserves the stationarity of subsets of \(\omega_1\).
    \item Proper forcing does not collapse \(\aleph_1\).
    \item The iteration of proper forcing notions with countable support is proper.
    \item Every forcing notion that is \(\sigma\)-closed, or ccc, or is a member of a wide class of other natural conditions, is proper.
\end{enumerate}
\end{theorem}

The third property --- the preservation of properness under countable support iterations --- is the most important. It provides a general framework for iterating forcing constructions while maintaining control over the resulting model. This result is the foundation for many of the most important independence proofs in set theory.

Shelah's theory of proper forcing is presented in his book \textit{Proper Forcing} (1982), which is the standard reference for the subject. The book develops the general theory and presents numerous applications, including consistency results about the continuum function, the structure of the real line, and problems in abelian group theory.

\subsection{Semiproper Forcing}

Later, Shelah extended the theory to \emph{semiproper forcing}, a broader class of forcing notions that includes proper forcing as a special case. Semiproper forcing was introduced to handle iterations where the forcing may not preserve \(\aleph_1\) but still maintains sufficient control.

\begin{definition}[Semiproper Forcing]
Let \(\mathbb{P}\) be a forcing notion. \(\mathbb{P}\) is \textbf{semiproper} if for every sufficiently large regular cardinal \(\theta\) and every countable elementary submodel \(N \prec H(\theta)\) with \(\mathbb{P} \in N\), and for every condition \(p \in \mathbb{P} \cap N\), there exists a condition \(q \leq p\) that is \(N\)-semigeneric: for every name \(\tau \in N\) for an ordinal, \(q \Vdash \tau \in N\).
\end{definition}

Shelah proved that the iteration of semiproper forcing notions with revised countable support (RCS) is semiproper, providing a powerful framework for forcing constructions that may involve the collapse of cardinals.

\subsection{The Shelah--Soifer Construction}

In a celebrated 1991 paper with Eli Soifer, Shelah used forcing to construct an unusual set of real numbers that does not have the Baire property and is also non-measurable, but whose construction does not require the axiom of choice in the full sense. This result illuminated the relationship between the axiom of choice, the existence of pathological sets of reals, and the role of large cardinal axioms.

\begin{theorem}[Shelah--Soifer, 1991]
It is provable in ZF + \(\DC\) (the axiom of dependent choice) that there exists a set of real numbers that does not have the Baire property. However, it is consistent with ZF + \(\DC\) that every set of reals is Lebesgue measurable.
\end{theorem}

This result showed that the existence of a non-measurable set of reals is not provable from ZF + \(\DC\) alone (it requires some additional choice, such as the existence of a well-ordering of \(\R\)), while the existence of a set without the Baire property is provable even with only dependent choice. This sharpened the understanding of the relationship between different regularity properties and the axiom of choice.

\section{J\'onsson Algebras}

A J\'onsson algebra is an algebra (a set with finitary operations) of cardinality \(\kappa\) that has no proper subalgebra of the same cardinality. The concept was introduced by Bjarni J\'onsson in the 1960s, and the problem of determining for which cardinals \(\kappa\) a J\'onsson algebra exists has deep connections with large cardinal axioms and the combinatorics of set theory.

\begin{definition}
An algebra \(\mathfrak{A} = (A, f_1, \dots, f_n)\) of cardinality \(\kappa\) is called a \textbf{J\'onsson algebra} if it has no proper subalgebra of cardinality \(\kappa\). Equivalently, if \(\mathfrak{A}\) is generated by every subset of \(A\) that has cardinality \(\kappa\).
\end{definition}

A \emph{J\'onsson cardinal} is a cardinal \(\kappa\) such that there is no J\'onsson algebra of cardinality \(\kappa\). The study of J\'onsson cardinals leads to the realm of large cardinals: the first J\'onsson cardinal (if it exists) is strongly inaccessible and is consistent with the existence of a measurable cardinal, but its exact strength is between that of an inaccessible and a measurable.

Shelah made fundamental contributions to the theory of J\'onsson algebras. He proved that:

\begin{itemize}
    \item For every successor cardinal \(\kappa^{+}\) that is not a J\'onsson cardinal, there exists a J\'onsson algebra of cardinality \(\kappa^{+}\).
    \item The first cardinal for which there is no J\'onsson algebra (if such a cardinal exists) must be a fixed point of the aleph function and strongly inaccessible.
    \item In the constructible universe \(L\), there are J\'onsson algebras at every successor cardinal, and the existence of a cardinal with no J\'onsson algebra implies certain large cardinal properties.
\end{itemize}

Shelah's work on J\'onsson algebras is closely connected with his work on the \emph{Keisler--Shelah theorem} and the structure of ultraproducts. The interplay between these different aspects of his work illustrates the remarkable coherence of his mathematical vision.

\section{Model Theory}

Beyond classification theory, Shelah has made numerous fundamental contributions to model theory understood broadly.

\subsection{Saturated Models}

A model \(M\) of a theory \(T\) is said to be \(\kappa\)-saturated if every type over a subset of \(M\) of size less than \(\kappa\) is realized in \(M\). Saturated models are important because they are the ``universal objects'' in the category of models: every model of size at most \(\kappa\) can be elementarily embedded into a \(\kappa^{+}\)-saturated model.

Shelah proved a deep characterization of the existence of saturated models in terms of the stability properties of the theory.

\begin{theorem}[Shelah's Saturated Model Theorem, 1971]
Let \(T\) be a complete countable first-order theory. Then:
\begin{itemize}
    \item If \(T\) is \(\omega\)-stable, then for every cardinal \(\kappa\), there exists a \(\kappa\)-saturated model of \(T\) of cardinality \(\kappa\).
    \item More generally, \(T\) has a \(\kappa\)-saturated model of cardinality \(\kappa\) for all sufficiently large \(\kappa\) if and only if \(T\) is stable.
    \item If \(T\) is unstable, then there are arbitrarily large cardinals \(\kappa\) for which no \(\kappa\)-saturated model of cardinality \(\kappa\) exists.
\end{itemize}
\end{theorem}

This theorem provides a striking connection between the existence of saturated models --- a purely model-theoretic concept --- and the stability classification.

\subsection{The Keisler--Shelah Theorem on Ultraproducts}

One of the most celebrated results in model theory is the Keisler--Shelah theorem, which characterizes elementary equivalence in terms of ultrapowers. The theorem states that two structures are elementarily equivalent if and only if they have isomorphic ultrapowers.

\begin{theorem}[Keisler--Shelah Theorem, 1967--1971]
Let \(L\) be a countable language, and let \(\mathfrak{A}\) and \(\mathfrak{B}\) be \(L\)-structures. Then \(\mathfrak{A} \equiv \mathfrak{B}\) (\(\mathfrak{A}\) and \(\mathfrak{B}\) are elementarily equivalent) if and only if there exists an ultrafilter \(D\) on some index set such that the ultrapowers \(\mathfrak{A}^I / D\) and \(\mathfrak{B}^I / D\) are isomorphic.

Moreover, for countable theories, the ultrafilter can be chosen to be a regular ultrafilter, and the theorem can be refined to relate the number of non-isomorphic ultrapowers to the stability properties of the theory.
\end{theorem}

The theorem was originally proved by Keisler in 1967 under the assumption of the GCH (the Generalized Continuum Hypothesis). Shelah removed this assumption in 1971, giving an unconditional proof using his theory of saturated models. The removal of the GCH hypothesis was considered a major breakthrough.

The Keisler--Shelah theorem has important consequences. It shows that elementary equivalence (the most basic notion of model theory) can be characterized purely in terms of ultrapowers, without any reference to the syntactic notion of satisfaction. It also provides a powerful tool for constructing saturated models and for analyzing the classification of theories.

\begin{corollary}
For any two countable elementarily equivalent structures \(\mathfrak{A}\) and \(\mathfrak{B}\), there exists an ultrafilter \(D\) such that \(\mathfrak{A}^{I}/D \cong \mathfrak{B}^{I}/D\). Thus the ultrapower construction can ``even out'' the differences between elementarily equivalent structures.
\end{corollary}

\subsection{Non-Structure Theory}

Shelah also developed the non-structure side of model theory, proving that in many contexts, the class of models of a theory cannot be classified by a reasonable set of invariants. The most famous of these results is the \emph{main gap theorem}, which classifies the number of non-isomorphic models of a first-order theory in each cardinality.

\begin{theorem}[Shelah's Main Gap Theorem, 1980s]
Let \(T\) be a complete countable first-order theory. For each uncountable cardinal \(\kappa\), let \(I(T,\kappa)\) denote the number of non-isomorphic models of \(T\) of cardinality \(\kappa\). Then either:
\begin{enumerate}
    \item \(I(T,\kappa) = 2^{\kappa}\) for all uncountable \(\kappa\) (the theory is \emph{wild}), or
    \item \(I(T,\kappa) < \beth_{\omega_1}(|\kappa|)\) for all uncountable \(\kappa\) (the theory is \emph{tame}), and there is a structure theory for the models of \(T\).
\end{enumerate}
Moreover, the dividing line is precisely whether the theory is classifiable (deep) according to Shelah's classification theory.
\end{theorem}

The main gap theorem is the culmination of Shelah's classification theory. It provides a complete answer to the question: how many models can a first-order theory have in each uncountable cardinality? The answer is either the maximal possible number (\(2^{\kappa}\)) or a relatively small number bounded by \(\beth_{\omega_1}(|\kappa|)\), and the two cases are distinguished by whether the theory is classifiable in Shelah's sense.

\section{Impact on Science and Technology}

\subsection{Mathematical Logic}

Shelah has been the dominant figure in mathematical logic for over four decades. His classification theory is the foundation of modern model theory, providing the framework in which all work on the structure of models is conducted. His pcf theory transformed cardinal arithmetic from a collection of isolated results into a systematic theory with deep theorems and open problems. His development of proper forcing and its iterations provided the technical foundation for much of modern set theory.

\subsection{Productivity and Style}

Shelah is one of the most prolific mathematicians of all time. With over 1200 refereed publications, his output exceeds that of nearly every other mathematician in history. His papers are known for their technical depth, their generality, and their tendency to contain far more than the advertised result. A typical Shelah paper might state one main theorem but prove dozens of subsidiary results that are themselves major contributions.

Shelah's writing style is distinctive. His papers are dense with notation and packed with ideas. He often publishes results in quick succession, refining and extending his previous work. His books --- \textit{Classification Theory} (1978, 2nd ed. 1990), \textit{Proper Forcing} (1982), \textit{Cardinal Arithmetic} (1994), and others --- are monumental works that have shaped entire fields.

\subsection{The Shelah School}

Shelah has supervised many PhD students, including Moti Gitik, Menachem Magidor, and others who have become leading figures in set theory and model theory. His influence extends far beyond his own students: every logician working in model theory, set theory, or the foundations of mathematics has been influenced by his work.

\subsection{Applications of His Work}

Shelah's work has found applications throughout mathematics. His classification theory has been applied in algebraic geometry (to the model theory of compact complex manifolds, differentially closed fields, and more), in number theory (to the model theory of valued fields and the study of Diophantine geometry), and in combinatorics (to the study of graphs and hypergraphs). His pcf theory has been used in combinatorial number theory and has led to new insights into the structure of ultrafilters and ultrapowers. His work on forcing is essential for understanding the consistency and independence of statements throughout mathematics.

\subsection{Open Problems}

Shelah's work has generated an enormous number of open problems. Some of the most important include:
\begin{enumerate}
    \item \textbf{The bound \(\aleph_{\omega_4}\):} Can Shelah's inequality \(\aleph_\omega^{\aleph_0} < \aleph_{\omega_4}\) be improved to \(\aleph_\omega^{\aleph_0} < \aleph_{\omega_3}\) or \(\aleph_{\omega_2}\)? The optimal bound is unknown.
    \item \textbf{The singular cardinals hypothesis (SCH):} Is it consistent that SCH fails at a cardinal of uncountable cofinality? Shelah's work has shown that such a failure is consistent relative to large cardinals, but the exact consistency strength is not known.
    \item \textbf{The classification of non-elementary classes:} Shelah has extended classification theory to various non-elementary frameworks (abstract elementary classes, \(\mathcal{L}_{\omega_1,\omega}\), etc.). The full classification of these classes remains an active area of research.
    \item \textbf{The continuum problem:} Despite Shelah's deep work on cardinal arithmetic, the continuum problem itself --- the question of the value of \(2^{\aleph_0}\) --- remains independent of ZFC, and most of Shelah's pcf bounds are about singular cardinals rather than \(\aleph_0\).
\end{enumerate}

\section*{Timeline}
\addcontentsline{toc}{section}{Timeline}\begin{tabular}{|p{2.5cm}|p{12cm}|}
\hline
\textbf{Year} & \textbf{Event} \\ \hline
1945 & Born on July 3 in Jerusalem, British Mandate for Palestine (now Israel) \\ \hline
1962 & Entered the Hebrew University of Jerusalem \\ \hline
1964 & B.Sc. in Mathematics and Physics, Hebrew University \\ \hline
1967 & M.Sc. in Mathematics, Hebrew University \\ \hline
1969 & Ph.D. under Michael Rabin at Hebrew University; dissertation on stable theories \\ \hline
1969--1970 & Postdoctoral research at the University of California, Berkeley \\ \hline
1970 & Joined the faculty of the Hebrew University of Jerusalem \\ \hline
1970s & Development of classification theory; the structure--nonstructure theorem; stability hierarchy \\ \hline
1974 & Independence of the Whitehead problem; beginning of his systematic work on forcing \\ \hline
1978 & Publication of \textit{Classification Theory}, the foundational text of the field \\ \hline
1978 & Introduction of proper forcing \\ \hline
1980s & Development of pcf (possible cofinalities) theory; Shelah's inequality \(\aleph_\omega^{\aleph_0} < \aleph_{\omega_4}\) \\ \hline
1982 & Publication of \textit{Proper Forcing} \\ \hline
1985--1990s & Development of semiproper forcing and revised countable support iterations \\ \hline
1991 & Shelah--Soifer construction: a set without the Baire property in ZF + DC \\ \hline
1994 & Publication of \textit{Cardinal Arithmetic}, systematizing pcf theory \\ \hline
1998 & Awarded the Israel Prize in Mathematics \\ \hline
2001 & Awarded the Wolf Prize in Mathematics (shared with Vladimir Arnold) \\ \hline
2011 & Awarded the EMET Prize in Mathematics \\ \hline
2015 & Over 1000 published papers (continues to publish actively) \\ \hline
\end{tabular}

\section*{Exercises for the Reader}
\addcontentsline{toc}{section}{Exercises}

\begin{enumerate}
    \item [I] Let \(T\) be a complete countable first-order theory that is \(\kappa\)-stable for some \(\kappa\). Show that \(T\) is \(\kappa\)-stable for all \(\kappa \geq 2^{\aleph_0}\). (This is a weaker version of the stability spectrum theorem.)
    \item [I] Prove that the theory of algebraically closed fields of characteristic 0 is \(\kappa\)-stable for all \(\kappa \geq \aleph_0\). Show that the theory of real closed fields is unstable.
    \item [I] Let \(A\) be a countable abelian group. Show that if \(\Ext^1(A,\Z) = 0\), then \(A\) is free. (This is Stein's theorem, the countable case of the Whitehead problem.)
    \item [I] Using the definition of forking, prove that in a stable theory, forking satisfies the extension property: if \(p\) is a type over \(A\) and \(A \subseteq B\), then there is a type \(q\) over \(B\) extending \(p\) that does not fork over \(A\).
    \item [I] Show that every \(\omega\)-stable theory is superstable, and every superstable theory is stable. Provide counterexamples showing that the converse implications fail.
    \item [A] Prove that if \(\mathfrak{a}\) is a set of regular cardinals with \(|\mathfrak{a}| = \aleph_0\) and \(\min(\mathfrak{a}) > \aleph_0\), then \(\max \pcf(\mathfrak{a}) \leq \aleph_{\omega_4}\). (This is the core of Shelah's inequality.)
    \item [A] Let \(V = L\). Show that every Whitehead group is free, using the \(\diamondsuit\) principle. (Hint: for a non-free group \(A\) of cardinality \(\aleph_1\), use \(\diamondsuit\) to construct a non-split extension of \(\Z\) by \(A\).)
    \item [A] Prove that if \(\MA + \neg\CH\) holds, then there exists a non-free Whitehead group of cardinality \(\aleph_1\). (Hint: construct a group \(A\) such that every countable subgroup of \(A\) is free.)
    \item [I] Let \(\mathbb{P}\) be a ccc forcing notion. Show that \(\mathbb{P}\) is proper. (This is one of the key examples of proper forcing.)
    \item [I] Show that \(\sigma\)-closed forcing is proper. Conclude that adding a Cohen subset to \(\omega_1\) is a proper forcing.
    \item [I] Prove that every model \(M\) of a stable theory \(T\) has a saturated elementary extension of cardinality \(|M|\). (This is the existence of saturated models in the stable case.)
    \item [A] Show that if two structures \(\mathfrak{A}\) and \(\mathfrak{B}\) are elementarily equivalent, then there is an ultrafilter \(D\) such that \(\mathfrak{A}^I/D\) and \(\mathfrak{B}^I/D\) are isomorphic. (This is the Keisler--Shelah theorem; you may use the GCH for the original Keisler proof.)
    \item [A] Let \(\kappa\) be a cardinal such that there is no J\'onsson algebra of cardinality \(\kappa\). Show that \(\kappa\) is strongly inaccessible. (This is a theorem of Shelah.)
    \item [I] Prove that the theory of the random graph is unstable. Show that it has the strict order property but not the independence property.
    \item [I] Let \(T\) be a superstable theory. Show that \(T\) has a \(\kappa\)-saturated model of cardinality \(\kappa\) for every regular cardinal \(\kappa\).
    \item [A] In the context of pcf theory, prove that if \(\mathfrak{a}\) is a progressive set of regular cardinals (i.e., \(|\mathfrak{a}| < \min(\mathfrak{a})\)), then \(\pcf(\mathfrak{a})\) has a maximum element.
    \item [I] Show that every ccc forcing that adds a real is proper but not \(\sigma\)-closed. Conclude that proper forcing is a strict generalization of both ccc and \(\sigma\)-closed forcing.
    \item [I] Let \(T\) be an \(\omega\)-stable theory. Show that the Morley rank of every complete type is defined and takes ordinal values. Prove that there exist types of arbitrarily high finite Morley rank.
    \item [A] Using the pcf theory, show that if \(\aleph_\omega\) is a strong limit cardinal, then \(2^{\aleph_\omega} < \aleph_{\omega_4}\). (Outline the key steps of Shelah's proof.)
    \item [A] Show that the class of models of a countable first-order theory \(T\) is \(\kappa\)-categorical for some \(\kappa\) (i.e., \(I(T,\kappa) = 1\)) if and only if \(T\) is \(\omega\)-stable and has no Vaughtian pairs. (This is the Morley--Shelah characterization of categoricity.)
\end{enumerate}

\section*{Glossary}
\addcontentsline{toc}{section}{Glossary}

\begin{description}
    \item[Classification theory] Shelah's systematic framework for analyzing first-order theories, dividing them into stable (tame) and unstable (wild) classes.
    \item[\(\kappa\)-stable] A theory is \(\kappa\)-stable if there are at most \(\kappa\) complete types over any model of size at most \(\kappa\).
    \item[Forking] Shelah's model-theoretic notion of independence; a formula forks over a set if it introduces a dependence that cannot be eliminated.
    \item[Structure--nonstructure theorem] Shelah's result that a first-order theory is either classifiable (admitting a structure theory) or wild (having the maximum number of models in all uncountable cardinalities).
    \item[Whitehead problem] The question of whether every abelian group \(A\) with \(\Ext^1(A,\Z) = 0\) is free; proved independent of ZFC by Shelah.
    \item[PCF theory] Shelah's possible cofinalities theory, a powerful framework for analyzing cardinal arithmetic on singular cardinals.
    \item[Shelah's inequality] \(\aleph_\omega^{\aleph_0} < \aleph_{\omega_4}\), the famous bound on the exponential of \(\aleph_\omega\).
    \item[\(\pcf(\mathfrak{a})\)] The set of possible cofinalities of a reduced product of regular cardinals; the central object of pcf theory.
    \item[\(\pp(\kappa)\)] The possible powers function, the supremum of certain cofinalities associated with \(\kappa\).
    \item[Proper forcing] A forcing notion that preserves stationarity of subsets of \(\omega_1\) and does not collapse \(\aleph_1\); introduced by Shelah.
    \item[Semiproper forcing] A generalization of proper forcing that allows for the collapse of \(\aleph_1\) while still maintaining control over iterations.
    \item[J\'onsson algebra] An algebra of cardinality \(\kappa\) with no proper subalgebra of cardinality \(\kappa\).
    \item[J\'onsson cardinal] A cardinal \(\kappa\) for which no J\'onsson algebra of cardinality \(\kappa\) exists; implies large cardinal properties.
    \item[Saturated model] A model realizing all types over subsets of size less than \(\kappa\).
    \item[Ultrapower] A construction \(\mathfrak{A}^I/D\) where \(\mathfrak{A}\) is a structure and \(D\) is an ultrafilter on an index set \(I\).
    \item[Keisler--Shelah theorem] Two structures are elementarily equivalent iff they have isomorphic ultrapowers.
    \item[Main gap theorem] Shelah's description of the possible numbers of non-isomorphic models of a first-order theory in each uncountable cardinality.
    \item[\(\diamondsuit\)] G\"odel's axiom of constructibility, asserting that every set is constructible; used by Shelah for the Whitehead problem.
    \item[\(\MA\)] Martin's Axiom, a combinatorial principle used by Shelah to construct non-free Whitehead groups.
    \item[Order property] A formula that defines an infinite linear order; its existence characterizes unstable theories.
    \item[Superstable] A theory that is stable and admits ordinal-valued ranks for all types.
    \item[\(\omega\)-stable] A theory that is \(\aleph_0\)-stable; equivalently, totally transcendental.
\end{description}

\section{Citations}

\subsection*{Primary Sources}
\begin{enumerate}
    \item Shelah, S. (1978). \textit{Classification Theory and the Number of Non-Isomorphic Models}. North-Holland.
    \item Shelah, S. (1975). The independence of the Whitehead problem. \textit{Journal of Algebra}, 35(1), 87--92.
    \item Shelah, S. (1974). A compactification for the space of models. \textit{Israel Journal of Mathematics}, 17(1), 1--15.
    \item Shelah, S. (1994). Cardinal arithmetic for singular cardinals. \textit{Annals of Pure and Applied Logic}, 69(3), 205--262.
\end{enumerate}

\subsection*{Biographies and Overviews}
\begin{enumerate}
    \item Hodges, W. (1999). A Shorter Shelah. \textit{Bulletin of Symbolic Logic}, 5(1), 1--22.
    \item Grossberg, R. (2002). Saharon Shelah's work in model theory. \textit{Notices of the American Mathematical Society}, 49(10), 1222--1229.
    \item American Mathematical Society. (2013). 2013 Steele Prize for Lifetime Achievement: Saharon Shelah. \textit{Notices of the American Mathematical Society}, 60(4), 472--475.
    \item Wolf Foundation. (2001). Saharon Shelah: Laureate in Mathematics. Retrieved from \url{https://www.wolffund.org.il/saharon-shelah/} (Accessed on October 26, 2023).
\end{enumerate}

\subsection*{Textbooks and Expository Works}
\begin{enumerate}
    \item Marker, D. (2002). \textit{Model Theory: An Introduction}. Springer.
    \item Jech, T. (2003). \textit{Set Theory} (3rd ed.). Springer.
    \item Tent, K., \& Ziegler, M. (2012). \textit{A Course in Model Theory}. Cambridge University Press.
    \item Kanamori, A. (2003). \textit{The Higher Infinite: Large Cardinals in Set Theory from Their Beginnings} (2nd ed.). Springer.
\end{enumerate}
